% !TeX program = xelatex
\documentclass[11pt,a4paper]{article}
\usepackage[utf8]{inputenc}
\usepackage[T1]{fontenc}
\usepackage[english,arabic]{babel}   % порядок важен: основной английский, дополнительный арабский
\usepackage{amsmath,amssymb,amsfonts,mathtools,amsthm}
\usepackage{geometry}
\geometry{left=1.25cm,right=1.25cm,top=1.25cm,bottom=3.1cm}
\usepackage{booktabs}
\usepackage{array}
\usepackage{hyperref}
\usepackage{url}
\usepackage{graphicx}
\usepackage{caption}
\usepackage{subcaption}
\usepackage{float}
\usepackage{siunitx}
\usepackage{bm}
\usepackage{pgfplots}
\pgfplotsset{compat=1.18}
\usepackage{xcolor}
\usepackage{titlesec}
\usepackage{fancyhdr}
\usepackage{microtype}
\usepackage{fontspec}
\usepackage{tcolorbox}
\usepackage{enumitem}

% Настройка шрифтов для латиницы (как в оригинале)
\setmainfont{Calibri}
\setsansfont{Segoe UI}[
BoldFont = Segoe UI Bold,
Ligatures = TeX
]

% Арабские шрифты – через babelfont для единообразия (включая курсив)
\babelfont[arabic]{rm}{Arial}[Scale=1.1, Script=Arabic]
\babelfont[arabic]{sf}{Arial}[Scale=1.1, Script=Arabic]
\babelfont[arabic]{tt}{Arial}[Scale=1.1, Script=Arabic]

% Цветовая схема YUCT
\definecolor{skyblue}{RGB}{0,102,204}
\definecolor{darkblue}{RGB}{0,0,139}
\definecolor{gold}{RGB}{255,215,0}

% Макросы YUCT
\newcommand{\Keff}{K_{\mathrm{eff}}}
\newcommand{\kappac}{\kappa_c}
\newcommand{\eps}{\varepsilon}
\newcommand{\betaf}{\beta}
\newcommand{\df}{d_{\!f}}
\newcommand{\Sodd}{S_{\mathrm{odd}}}
\newcommand{\Seven}{S_{\mathrm{even}}}
\newcommand{\Mpl}{M_{\mathrm{Pl}}}
\newcommand{\Lpl}{l_{\mathrm{Pl}}}
\newcommand{\tpl}{t_{\mathrm{Pl}}}
\newcommand{\Lmin}{\ell_{\min}}
\newcommand{\LUV}{\Lambda_{\mathrm{UV}}}

% Теоремы
\newtheorem{theorem}{نظرية}[section]
\newtheorem{proposition}[theorem]{اقتراح}
\newtheorem{definition}[theorem]{تعريف}
\theoremstyle{remark}
\newtheorem{remark}[theorem]{ملاحظة}

% Форматирование секций (выравнивание заголовков вправо для RTL)
\titleformat{\section}{\LARGE\bfseries\color{darkblue}\raggedleft}{\thesection}{1em}{}
\titleformat{\subsection}{\Large\bfseries\color{darkblue}\raggedleft}{\thesubsection}{1em}{}
\titleformat{\subsubsection}{\normalsize\bfseries\color{darkblue}\raggedleft}{\thesubsubsection}{1em}{}

% Верхний и нижний колонтитулы
\fancypagestyle{firstpage}{%
	\fancyhf{}%
	\renewcommand{\headrulewidth}{-5pt}%
	\renewcommand{\footrulewidth}{-5pt}%
	\fancyfoot[C]{%
		\begin{minipage}[t]{\textwidth}
			\centering
			\textcolor{gold}{\rule{\textwidth}{1.6pt}}\\[5pt]
			{\small \textsuperscript{1}Yakushev Research, \url{https://yuct.org/}} \hfill
			{\small بروتوكول YPSDC: \url{https://ypsdc.com/}}\\[-2pt]
			\textcolor{gold}{\rule{\textwidth}{1.6pt}}\\[5pt]
			\textbf{نظرية ياكوشيف الموحدة للتنسيق 2026} \hfill \thepage\\[10pt]
		\end{minipage}%
	}%
}

\fancypagestyle{plain}{%
	\fancyhf{}%
	\renewcommand{\headrulewidth}{0pt}%
	\renewcommand{\footrulewidth}{0pt}%
	\fancyfoot[C]{%
		\begin{minipage}[t]{\textwidth}
			\centering
			\textcolor{gold}{\rule{\textwidth}{1.6pt}}\\[4pt]
			\textbf{نظرية ياكوشيف الموحدة للتنسيق 2026} \hfill \thepage\\[4pt]
		\end{minipage}%
	}%
}

\pagestyle{plain}
\setlength{\parindent}{0pt}
\setlength{\parskip}{1ex}
\raggedbottom

\hypersetup{
	colorlinks=true,
	linkcolor=blue,
	citecolor=blue,
	urlcolor=blue,
}

% Шапка документа
\newcommand{\makeheader}{%
	\noindent
	\begin{minipage}{0.17\textwidth}
		\raggedright
		{\fontspec{Segoe UI}[BoldFont=Segoe UI Bold]\bfseries\color{skyblue}\fontsize{46}{46}\selectfont YUCT}%
	\end{minipage}%
	\hspace{30pt}
	\begin{minipage}{0.32\textwidth}
		\raggedright
		{\sffamily\fontsize{11}{13}\selectfont YAKUSHEV\\ UNIFIED\\ COORDINATION THEORY}%
	\end{minipage}%
	\hfill
	\begin{minipage}{0.38\textwidth}
		\raggedleft
		{\sffamily\bfseries الملحق UVD}\\ 
		{\sffamily الإصدار 1.0 / حزيران 2026}\\ 
		{\sffamily\footnotesize \url{https://doi.org/10.5281/zenodo.18444598}}%
	\end{minipage}
	\par\vspace{5pt}
	\noindent\textcolor{gold}{\rule{\textwidth}{1.6pt}}
	\par\vspace{0pt}
}

\begin{document}
	\thispagestyle{firstpage}
	\makeheader
	
	\begin{otherlanguage}{arabic}
		\vspace{0.5cm}
		{\LARGE\bfseries الملحق UVD: حل الانفراجات فوق البنفسجية في YUCT \\
			\Large نظرية الحقل الكمي المنتهية من عامل شكل التنسيق والمقياس الفيزيائي لشبكة التنسيق \par}
		\vspace{0.3cm}
		
		{\large Alexey V. Yakushev\textsuperscript{1}} \\
		\vspace{0.2cm}
		
		{\color{darkblue}\textbf{\LARGE الملخص}} \par
		{\large
			\noindent
			تقترح نظرية ياكوشيف الموحدة للتنسيق (YUCT) أن الانفراجات فوق البنفسجية في نظرية الحقل الكمي يتم التخلص منها بواسطة عامل شكل التنسيق \(F(q^2) = \exp[-(q^2/\LUV^2)^{2/3}]\) الناشئ عن السعة المعلوماتية المنتهية للفراغ. في الصياغة الأصلية، تم تحديد المقياس فوق البنفسجي \(\LUV\) بكتلة بلانك، \(\LUV = 3\pi\Mpl \approx 1.15\times 10^{20}\) GeV، مما يؤدي إلى تصحيح إشعاعي منتهٍ هائل لكتلة هيغز. نصّح خطأً حسابياً في الثابت \(C_0\) الذي يظهر في تصحيح الكتلة السلمية: القيمة الصحيحة هي \(C_0 = 3\sqrt{\pi}/(8\sqrt{2}) \approx 0.46999\)، لتحل محل القيمة المبلغ عنها سابقاً \(C_0 \approx 0.626\). وبهذا التصحيح، يصبح المساهمة الإشعاعية لكتلة هيغز من شبكة بمقياس بلانك \(\delta m \approx 1.57\times 10^{18}\) GeV، مما يعيد إشكالية الضبط الدقيق الحاد ما لم يتم تعويض الكتلة العارية بدقة. نفحص بديلين منطقيين: (i) صيغة الكتلة في YUCT \(m = \Mpl(2/3)^L\xi_f\) تحدد الكتلة الفيزيائية بالفعل، ويتم امتصاص التصحيح الإشعاعي بواسطة حد مضاد منتهٍ – حل متسق تقنياً لكنه غير مرضي جمالياً؛ (ii) لا تحتوي لاغرانجية الأساسية على مقياس بلانك على حد كتلة عارية (ثبات مقياسي)، ويتم توليد كتلة هيغز بأكملها بواسطة التصحيح الإشعاعي. في هذا السيناريو الثاني، تستلزم الاتساق الذاتي أن يكون المقياس فوق البنفسجي \(\LUV \approx 8.96\) TeV، وهي قيمة يمكن الوصول إليها في مصادم الهدرونات الكبير (LHC). عندئذٍ يتنبأ عامل شكل التنسيق بانحرافات ملحوظة في المقاطع العرضية للاستطارة عالية الطاقة عند مقياس TeV. نقدم كلا السيناريوهين، وتبريرهما الرياضي، وعواقبهما التجريبية بهدف إخضاع YUCT لاختبار تجريبي صارم.
		}
		
		\vspace{0.2cm}
		\noindent
		{\color{darkblue}\textbf{الكلمات المفتاحية:}} YUCT، الانفراجات فوق البنفسجية، عامل شكل التنسيق، مشكلة التسلسل الهرمي، كتلة هيغز، الثبات المقياسي، نظرية الحقل الكمي غير الموضعية، فينومينولوجيا LHC.
		\vspace{0.1cm}
		\noindent
		{\color{darkblue}\textbf{الاستشهاد:}} Yakushev AV. الملحق UVD: حل الانفراجات فوق البنفسجية في YUCT. 2026. DOI: \url{https://doi.org/10.5281/zenodo.18444598} (هذا المجلد).
		
		\vspace{0.5cm}
		
		\section{مقدمة}
		\label{sec:intro}
		
		تقدم نظرية ياكوشيف الموحدة للتنسيق (YUCT) حلاً أنيقاً للانفراجات فوق البنفسجية في نظرية الحقل الكمي (QFT). بدلاً من إدخال قطع زخم اعتباطي أو إعادة تطبيع لانهائي، تفترض YUCT أن الفراغ الفيزيائي يمتلك "عرض نطاق" منتهياً لنقل المعلومات الكمومية. يتم تحديد هذا العرض بواسطة \textbf{عامل شكل التنسيق} \(F(q^2)\) الذي يكبح الزخم الافتراضي فوق مقياس مميز \(\LUV\) بشكل أسي. يتم تحديد البنية الرياضية لهذا العامل بقانون الخطأ العام في YUCT، \(\varepsilon \propto K_{\mathrm{eff}}^{-\beta}\) حيث \(\beta = 2/3\).
		
		في الإصدارات السابقة من هذا الملحق، طابقنا \(\LUV\) مع كتلة بلانك، \(\LUV = 3\pi\Mpl \approx 1.15\times 10^{20}\) GeV – وهو مقياس مشتق من الطول الأدنى \(\Lmin = \frac{2}{3}\Lpl\) الذي يمليه الحلقة الجبرية الأولية لـ YUCT. مع هذه المطابقة، تصبح جميع تكاملات الحلقات منتهية بشكل واضح: يتم استبدال الانفراجات اللوغاريتمية بـ \(\ln(\LUV/m)\)، وتختفي انفراجات القوى تماماً.
		
		مع ذلك، تكشف إعادة حساب دقيقة لتصحيح الكتلة السلمية عن خطأ حسابي في الثابت \(C_0\) الذي يتحكم في حجم البقية المنتهية. القيمة الصحيحة لـ \(C_0\) أصغر مما كان مبلغاً عنه سابقاً، مما يقلل التصحيح الإشعاعي ولكنه لا يزال يتركه أكبر بعدة مراتب من كتلة هيغز المرصودة. هذا يجبرنا على إعادة النظر نقدياً في المقياس الفيزيائي \(\LUV\).
		
		لهذا الملحق هدفان. أولاً، نقدم الاشتقاقات الرياضية المصححة بالتفصيل الكامل، لضمان دقة جميع التعبيرات وقابليتها للتكرار. ثانياً، نستكشف الآثار الفيزيائية للصيغ المصححة، بما في ذلك اقتراح جذري: المقياس الأساسي لشبكة التنسيق قد لا يكون مقياس بلانك على الإطلاق، بل مقياساً أقل بكثير، ربما يكون في متناول تجارب المصادمات الحالية.
		
		\section{عامل شكل التنسيق والمقياس فوق البنفسجي}
		\label{sec:formfactor}
		
		عامل شكل التنسيق في YUCT هو
		\begin{equation}
			F(q^2) = \exp\!\Bigl[-\Bigl(\frac{q^2}{\LUV^2}\Bigr)^{\!\beta}\,\Bigr], \qquad \beta = \frac{2}{3}.
			\label{eq:F}
		\end{equation}
		يتم ضرب جميع مكبرات النشر في النموذج المعياري بـ \(F(p^2/\LUV^2)\). يتم اشتقاق المقياس \(\LUV\) أصلاً من الطول الأدنى \(\Lmin = \frac{2}{3}\Lpl\) عبر علاقة كومبتون \(\LUV = 2\pi\hbar c/\Lmin = 3\pi\Mpl\).
		
		هذا العامل يجعل جميع تكاملات الحلقات متقاربة تماماً في المنطقة فوق البنفسجية. البرهان بسيط: لأي متعددة حدود \(P(k)\) في زخم الحلقة، فإن الجداء \(P(k) \cdot \prod_i F(k_i^2/\LUV^2)\) يتناقص أسرع من أي قوة عكسية لـ \(|k|\) عندما \(|k| \to \infty\)، لأن الدالة الأسية هي المسيطرة.
		
		\section{طاقة الإلكترون الذاتية: حل انفراج لوغاريتمي}
		\label{sec:electron}
		
		طاقة الإلكترون الذاتية عند حلقة واحدة في الديناميكا الكهربائية الكمية (QED) هي من النوع المتباعد لوغاريتمياً في QTF التقليدية. مع عامل شكل YUCT، يصبح التكامل الإقليدي
		\begin{equation}
			\Sigma_{\mathrm{YUCT}} \propto \int_0^{\infty} \frac{dk_E}{k_E}\; \exp\!\Bigl[-2\Bigl(\frac{k_E^2}{\LUV^2}\Bigr)^{\!2/3}\,\Bigr].
			\label{eq:elec_int}
		\end{equation}
		بالتعويض \(t = (k_E^2/\LUV^2)^{2/3}\) نحصل على \(dk_E/k_E = \frac{3}{4} dt/t\)، ويصبح التكامل
		\begin{equation}
			I_{\mathrm{log}} = \frac{3}{4} \int_{t_{\min}}^{\infty} \frac{dt}{t}\, e^{-2t} \approx \ln\frac{\LUV}{m_e},
			\label{eq:elec_result}
		\end{equation}
		حيث \(t_{\min} \sim (m_e^2/\LUV^2)^{2/3}\). يتم استبدال الانفراج اللوغاريتمي باللوغاريتم المنتهي \(\ln(\LUV/m_e) \approx 53.8\). الحساب أولي ولا يحتاج إلى تصحيح.
		
		\section{تصحيح الكتلة السلمية: تصحيح الثابت \(C_0\)}
		\label{sec:scalar}
		
		\subsection{التكامل وحسابه}
		تصحيح الكتلة السلمية عند حلقة واحدة في نظرية \(\phi^4\)، بعد استدارة ويك وإدخال عامل الشكل، هو
		\begin{equation}
			\delta m^2_{\mathrm{YUCT}} = \frac{\lambda}{16\pi^2} \int_0^{\infty} \frac{k_E^3\, dk_E}{k_E^2 + m^2}\; F^2\!\Bigl(\frac{k_E^2}{\LUV^2}\Bigr), \qquad F^2(x) = e^{-2x^{2/3}}.
			\label{eq:scalar_YUCT}
		\end{equation}
		
		بإدخال المتغير اللامبعد \(t = k_E^2/\LUV^2\) وبالنشر لـ \(m \ll \LUV\)، نحصل على الحد الرئيسي
		\begin{equation}
			\delta m^2_{\mathrm{YUCT}} = \frac{\lambda}{32\pi^2}\,\LUV^2 \int_0^{\infty} e^{-2t^{2/3}} dt + \mathcal{O}(m^2/\LUV^2).
			\label{eq:scalar_lead}
		\end{equation}
		
		سابقاً (الإصدار 3.1) قُدِّر التكامل \(C_0 \equiv \int_0^{\infty} e^{-2t^{2/3}} dt \approx 0.626\). هذه القيمة خاطئة. نحسبه الآن بدقة.
		
		\subsection{الحساب الدقيق لـ \(C_0\)}
		\begin{theorem}[القيمة الدقيقة لـ \(C_0\)]
			\label{thm:C0}
			\[
			C_0 = \int_0^{\infty} e^{-2t^{2/3}} dt = \frac{3\sqrt{\pi}}{8\sqrt{2}} \approx 0.46999.
			\]
		\end{theorem}
		
		\begin{proof}
			نضع \(u = 2t^{2/3}\). إذن \(t = (u/2)^{3/2}\) و
			\[
			dt = \frac{3}{2} (u/2)^{1/2} \cdot \frac{1}{2} du = \frac{3}{4\sqrt{2}} u^{1/2} du.
			\]
			بالتالي
			\[
			C_0 = \int_0^{\infty} e^{-u} \cdot \frac{3}{4\sqrt{2}} u^{1/2} du = \frac{3}{4\sqrt{2}} \int_0^{\infty} u^{1/2} e^{-u} du = \frac{3}{4\sqrt{2}} \, \Gamma(3/2).
			\]
			باستخدام \(\Gamma(3/2) = \frac{\sqrt{\pi}}{2}\)، نحصل على
			\[
			C_0 = \frac{3}{4\sqrt{2}} \cdot \frac{\sqrt{\pi}}{2} = \frac{3\sqrt{\pi}}{8\sqrt{2}} \approx 0.46999.
			\]
		\end{proof}
		
		\subsection{تصحيح الكتلة السلمية المصحح}
		باستخدام القيمة الصحيحة لـ \(C_0\)، يصبح الحد الرئيسي لتصحيح الكتلة السلمية
		\begin{equation}
			\boxed{\delta m^2_{\mathrm{YUCT}} = \frac{\lambda}{32\pi^2}\,\LUV^2 \cdot \frac{3\sqrt{\pi}}{8\sqrt{2}} + \mathcal{O}(m^2/\LUV^2)}.
			\label{eq:scalar_corrected}
		\end{equation}
		
		بالنسبة لبوزون هيغز، \(\lambda = m_H^2/(2v^2) \approx 0.13\) (حيث \(v \approx 246\) GeV). باستخدام مطابقة مقياس بلانك \(\LUV = 3\pi\Mpl \approx 1.15\times 10^{20}\) GeV:
		\begin{align}
			\delta m^2_H &\approx \frac{0.13}{32\pi^2} \cdot (1.15\times 10^{20})^2 \cdot 0.46999 \notag\\
			&\approx 2.47 \times 10^{36}\ \mathrm{GeV}^2, \notag\\
			\delta m_H &\approx 1.57 \times 10^{18}\ \mathrm{GeV}.
			\label{eq:planck_correction}
		\end{align}
		
		هذه القيمة تتجاوز كتلة هيغز المرصودة (\(m_H^{\mathrm{exp}} \approx 125\) GeV) بعامل \(1.26 \times 10^{16}\). إذا كانت الكتلة العارية في لاغرانجية صفراً، فإن الكتلة الفيزيائية تتحدد بالكامل بهذا التصحيح الإشعاعي، ومن الواضح أن القيمة المتوقعة غير متوافقة مع الرصد.
		
		\section{الآثار الفيزيائية: سيناريوهان}
		\label{sec:scenarios}
		
		الحساب المصحح يجبرنا على مواجهة مشكلة التسلسل الهرمي مباشرة. نفحص بديلين متسقين منطقياً.
		
		\subsection{السيناريو A: شبكة بمقياس بلانك مع حدود مضادة منتهية}
		في هذا السيناريو، يحتفظ YUCT بمطابقة \(\LUV = 3\pi\Mpl\). تعطى الكتلة الفيزيائية لأي جسيم بواسطة صيغة كتلة YUCT
		\begin{equation}
			m_{\mathrm{phys}} = \Mpl \left(\frac{2}{3}\right)^L \xi_f,
			\label{eq:mass_formula}
		\end{equation}
		حيث \(L\) هو عمق التنسيق و \(\xi_f\) عامل رنين (انظر الملحق J). الكتلة العارية في لاغرانجية هي معامل حر، يُختار ليلغي التصحيح الإشعاعي ويعيد إنتاج الكتلة المرصودة.
		
		رياضياً، هذا متسق تماماً: النظرية منتهية، والحد المضاد المطلوب منته أيضاً، ولا نواجه أبداً ما لا نهاية. صيغة كتلة YUCT، المعادلة (\ref{eq:mass_formula})، تقدم شرط إعادة التطبيع الذي يثبت الحد المضاد.
		
		فيزيائياً، مع ذلك، يتطلب هذا السيناريو أن تلغي الكتلة العارية والتصحيح الإشعاعي بعضهما بدقة جزء واحد في \(10^{16}\) – ضبط دقيق من النوع الذي كان من المفترض أن تتجنبه مشكلة التسلسل الهرمي. YUCT لا يفسر لماذا يجب اختيار الكتلة العارية لتحقيق هذا الإلغاء؛ إنه يستبدل ببساطة الطرح اللانهائي المعتاد بطرح منتهٍ هائل.
		
		\subsection{السيناريو B: أصل ثابت مقياسياً وشبكة بمقياس TeV}
		بديل أكثر قبولاً من الناحية الجمالية هو التخلي عن مطابقة \(\LUV\) لمقياس بلانك. إذا كانت النظرية الأساسية عند مقياس بلانك لا تحتوي على حدود كتلة عارية (أي أنها ثابتة مقياسياً)، فإن جميع الكتل تتولد إشعاعياً بواسطة عامل شكل التنسيق. في هذه الحالة، كتلة هيغز الفيزيائية \emph{هي} التصحيح الإشعاعي:
		\begin{equation}
			m_H^2 = \delta m^2_{\mathrm{YUCT}} = \frac{\lambda}{32\pi^2}\,\LUV^2 \cdot C_0.
			\label{eq:scale_inv}
		\end{equation}
		
		بحل المعادلة للحصول على \(\LUV\):
		\begin{equation}
			\LUV = \sqrt{\frac{32\pi^2\, m_H^2}{\lambda\, C_0}}.
			\label{eq:LUV_solve}
		\end{equation}
		
		بتعويض \(C_0 = 3\sqrt{\pi}/(8\sqrt{2}) \approx 0.46999\)، \(\lambda = 0.13\)، و \(m_H = 125\) GeV، نحصل على
		\begin{equation}
			\boxed{\LUV \approx 8.96\ \mathrm{TeV}}.
			\label{eq:LUV_TeV}
		\end{equation}
		
		هذه القيمة هي تنبؤ للنظرية، وليست معاملاً حراً. إنها تضع المقياس الأساسي للتنسيق في نطاق TeV، مما يجعله في متناول التجارب.
		
		\subsection{التفسير الفيزيائي للسيناريو B}
		إذا كان \(\LUV \approx 9\) TeV، فإن شبكة التنسيق "تتبلور" عند المقياس الكهروضعيف، وليس عند مقياس بلانك. يبدأ عامل الشكل \(F(q^2) = \exp[-(q^2/\LUV^2)^{2/3}]\) في كبح الزخم الافتراضي فوق بضعة TeV. لهذا عواقب ملحوظة فورية:
		
		\begin{enumerate}[label=(\arabic*)]
			\item \textbf{المقاطع العرضية لدريل-يان والنفاثات المزدوجة} في LHC يجب أن تُظهر عجزاً في الأحداث ذات الكتل الثابتة العالية مقارنة بتوقعات النموذج المعياري. يصبح الكبح ملحوظاً لـ \(m_{jj} \gtrsim 5\) TeV.
			\item \textbf{إنتاج أزواج هيغز} وتشتت بوزونات المتجه يجب أن يُظهرا اعتماداً شاذاً على الطاقة لأن عامل الشكل يعدل اقتران هيغز الذاتي.
			\item \textbf{انحرافات في المراصد الكهروضعيفة الدقيقة} قد تكون متوقعة، لكنها على الأرجح صغيرة لأن قياسات LEP تستكشف مقاييس أقل بكثير من \(\LUV\).
		\end{enumerate}
		
		الأهم هو أن الطول الأدنى في هذا السيناريو هو \(\Lmin = 2\pi\hbar c/\LUV \approx 1.38\times 10^{-19}\) m، وهو أكبر بكثير من طول بلانك (\(\sim 10^{-35}\) m). هذا يعني أن خشونة الزمكان ليست تأثيراً جاذبياً كمياً بل ظاهرة بمقياس كهروضعيف، مرتبطة على الأرجح بآلية هيغز نفسها.
		
		\subsection{الجسر الكسيري بين المقياس الكهروضعيف والثابت الكوني}
		\label{subsec:bridge}
		إذا كان السيناريو B صحيحاً والمقياس الأساسي للتنسيق هو \(\LUV \approx 8.96\) TeV، يظهر ارتباط عميق بعلم الكون. في YUCT، الفراغ ليس وسطاً متصلاً واحداً بل تسلسلاً هرمياً كسيرياً من طبقات التنسيق. يصف المقياس \(\LUV\) الطبقة التي تشهد فيها حقول النموذج المعياري عامل شكل التنسيق. ومع ذلك، توجد طبقات أعمق (تحت حمراء) تقابل مقاييس طاقة أقل تدريجياً. الثابت الكوني، أي كثافة الطاقة المظلمة المرصودة، يتولد بواسطة أعمق طبقة تحت حمراء، التي يُحدد مقياسها \(\Lambda_{\mathrm{IR}}\) بشرط أن طاقة نقطة الصفر لجميع حقول الكم، عند تكاملها عبر التسلسل الهرمي الكسيري الكامل، تعيد إنتاج القيمة المرصودة \(\rho_{\Lambda} \approx 2.51 \times 10^{-47}\) GeV\(^4\) (انظر الملحق M).
		
		العلاقة بين المقياس فوق البنفسجي \(\LUV\) والمقياس تحت الأحمر \(\Lambda_{\mathrm{IR}}\) يحكمها الكم الأساسي للقياس في YUCT:
		\begin{equation}
			q = \left(\frac{3}{2}\right)^{1/3} \approx 1.1447.
		\end{equation}
		يفصل بين المقياسين عدد صحيح \(N\) من طبقات التنسيق:
		\begin{equation}
			\LUV = \Lambda_{\mathrm{IR}} \cdot q^{N}.
			\label{eq:bridge}
		\end{equation}
		
		من الثابت الكوني المرصود نستخلص \(\Lambda_{\mathrm{IR}} \approx 0.0032\) eV (انظر الملحق M، Eq.~(M.47)). وباستخدام \(\LUV \approx 8.96 \times 10^{12}\) eV، نحصل على
		\begin{equation}
			N = \frac{\ln(\LUV / \Lambda_{\mathrm{IR}})}{\ln q}
			= \frac{\ln(8.96 \times 10^{12} / 0.0032)}{\frac{1}{3}\ln(3/2)}
			\approx \frac{35.564}{0.135155} \approx 263.1.
			\label{eq:N_layers}
		\end{equation}
		
		بتقريب إلى أقرب نصف صحيح (كما تطلبه الحلقة الجبرية لـ YUCT للطبقات الفرميونية)، نحصل على العدد البنيوي الأساسي
		\begin{equation}
			\boxed{N = 263.0}.
		\end{equation}
		
		هذه نتيجة رائعة. المقياس الكهروضعيف (كتلة بوزون هيغز) والثابت الكوني مرتبطان بـ 263 خطوة متقطعة بالضبط من التسلسل الهرمي للتنسيق، بدون أي معاملات حرة. صغر الثابت الكوني المرصود ليس لغزاً؛ إنه يعكس ببساطة العمق الهائل لشبكة التنسيق التي تفصل فيزياء الجسيمات عن علم الكون.
		
		\textbf{تفسير:} كل طبقة من التسلسل الهرمي تعمل كـ"محول مقياس": طاقة نقطة الصفر لطبقة معينة يتم كبحها بواسطة عامل شكل التنسيق للطبقة التالية الأعمق. طاقة الفراغ الكلية هي مجموع مساهمات جميع الطبقات، كل منها أصغر من سابقتها بعامل \(\beta = 2/3\). تتقارب المتسلسلة الهندسية إلى قيمة منتهية، وهي بالضبط كثافة الطاقة المظلمة المرصودة (انظر الملحق M). العدد الصحيح \(N \approx 263\) هو عدد الطبقات التي تناسب بين المقياس الكهروضعيف ومقياس كتلة النيوترينو (أعمق طبقة يمكن الوصول إليها فيزيائياً). أن هذا العدد يظهر كنصف صحيح هو فحص اتساق غير بديهي لإطار YUCT، مما يدعم السيناريو B بشكل إضافي.
		
		\subsection{مقارنة بين السيناريوين}
		\begin{table}[H]
			\centering
			\caption{مقارنة السيناريو A والسيناريو B}
			\label{tab:scenarios}
			\small
			\begin{tabular}{@{}lcc@{}}
				\toprule
				\textbf{الخاصية} & \textbf{السيناريو A} & \textbf{السيناريو B} \\
				\midrule
				المقياس فوق البنفسجي \(\LUV\) & \(3\pi\Mpl \approx 10^{20}\) GeV & \(\approx 8.96\) TeV \\
				الطول الأدنى \(\Lmin\) & \(\frac{2}{3}\Lpl \approx 10^{-35}\) m & \(\approx 10^{-19}\) m \\
				الكتلة العارية في لاغرانجية & غير صفرية (ضبط دقيق) & صفر (ثبات مقياسي) \\
				التصحيح الإشعاعي \(\delta m_H\) & \(10^{18}\) GeV & 125 GeV \\
				إمكانية الوصول التجريبي & مقياس بلانك (غير ممكن) & طاقات LHC (ممكنة) \\
				ضبط دقيق مطلوب & نعم (جزء واحد في \(10^{16}\)) & لا \\
				\bottomrule
			\end{tabular}
		\end{table}
		
		\section{تنبؤات قابلة للتزييف للسيناريو B}
		\label{sec:predictions}
		
		إذا كان السيناريو B صحيحاً، يمكن اختبار التنبؤات الكمية التالية في LHC والمصادمات المستقبلية:
		
		\begin{enumerate}[label=(\arabic*)]
			\item \textbf{كبح مقطع دريل-يان.} نسبة مقطع دريل-يان المقاس إلى مقطع النموذج المعياري كدالة في الكتلة الثابتة للبتونين \(m_{\ell\ell}\) يجب أن تتبع
			\[
			R(m_{\ell\ell}) \approx \exp\!\Bigl[-2\Bigl(\frac{m_{\ell\ell}^2}{\LUV^2}\Bigr)^{\!2/3}\,\Bigr],
			\]
			مع \(\LUV \approx 8.96\) TeV. يُتوقع كبح بنسبة 10\% عند \(m_{\ell\ell} \approx 5.4\) TeV، ويتزايد إلى عامل \(\sim 2\) عند \(m_{\ell\ell} \approx 10\) TeV.
			
			\item \textbf{إنتاج أزواج هيغز.} يعدل عامل الشكل اقتران هيغز الذاتي عند انتقال زخم عال. يجب أن ينحرف توزيع الكتلة الثابتة لثنائي هيغز عن النموذج المعياري لـ \(m_{HH} \gtrsim 3\) TeV.
			
			\item \textbf{كتل الفرميونات.} إذا كانت جميع كتل الفرميونات تتولد أيضاً إشعاعياً، فيجب اختيار قرانات يوكاوا بحيث تعيد التصحيحات الإشعاعية إنتاج الطيف الكتلي المرصود. هذا يفرض شرط اتساق ذاتي يربط أعماق التنسيق \(L_f\) وعوامل الرنين \(\xi_f\) بالتصحيحات الإشعاعية. تحليل مفصل يُؤجل إلى عمل مستقبلي.
		\end{enumerate}
		
		\subsection{ملاحظة حول ثابت البنية الدقيقة}
		اشتقاق \(\alpha^{-1} \approx 137.036\) من الثابت الطوبولوجي \(N_{\mathrm{gauge}} = 12\) (الملحق G) مستقل عن قيمة \(\LUV\). قرانات المقياس تجري بشكل لوغاريتمي؛ قيمها الابتدائية عند \(\LUV\) تُثبّت بطوبولوجيا \(F_{18}\). في السيناريو A، \(\LUV \sim \Mpl\)، والجري اللوغاريتمي من مقياس بلانك إلى المقياس الكهروضعيف يعطي \(\alpha^{-1} \approx 137.036\). في السيناريو B، \(\LUV \sim 9\) TeV، والجري يكون على فترة أقصر بكثير. عندئذٍ ينطبق التنبؤ الطوبولوجي لـ \(\alpha^{-1}\) على مقياس TeV، ويتم الحصول على القيمة منخفضة الطاقة بتطور لوغاريتمي صغير. هذا سيغير التنبؤ العددي لـ \(\alpha^{-1}\) ويجب التحقق منه مع قيمة CODATA. تشير التقديرات الأولية إلى أن الانزياح ضمن عدم اليقين الحالي للتنبؤ الطوبولوجي، لكن حساباً دقيقاً مطلوب وسيُقدّم في مراجعة مستقبلية للملحق G.
		
		\section{مناقشة: هل YUCT نظرية كل شيء أم نظرية المقياس الكهروضعيف؟}
		\label{sec:discussion}
		
		إطار YUCT الأصلي طمح إلى تفسير جميع المقاييس الفيزيائية من كتلة بلانك إلى كتلة الإلكترون عبر تسلسل هرمي واحد للتنسيق. يكشف حساب الكتلة السلمية المصحح عن توتر داخل هذا الطموح: إذا كانت شبكة التنسيق تعمل على مقياس بلانك، فإن التصحيحات الإشعاعية هائلة ويجب إلغاؤها بكتل عارية مضبوطة بدقة.
		
		السيناريو B يحل هذا التوتر بإعادة تفسير سلم الكتلة في YUCT. بدلاً من إنتاج الكتل مباشرة من خلال العمق \(L\)، يحدد السلم \textbf{قرانات يوكاوا} التي بدورها تولد الكتل إشعاعياً على المقياس \(\LUV \sim 9\) TeV. يصبح مقياس بلانك عندئذٍ هو المقياس الذي تبدأ فيه شبكة التنسيق \textit{بالتشكل}، لكن الخشونة الفيزيائية – المقياس الذي يصبح فيه عامل الشكل مهماً – أقل بكثير.
		
		إعادة التفسير هذه تحافظ على الجمال الجبري لـ YUCT مع إزالة مشكلة الضبط الدقيق. كما أنها تحول YUCT من نظرية تخمينية للجاذبية الكمية إلى نموذج قابل للتزييف لفيزياء الكهروضعيف، يمكن اختباره في المصادمات الحالية والمستقبلية.
		
		\section{الاستنتاج}
		\label{sec:conclusion}
		
		صححنا خطأ حسابياً في تصحيح الكتلة السلمية داخل YUCT، وحصلنا على القيمة التحليلية الدقيقة \(C_0 = 3\sqrt{\pi}/(8\sqrt{2}) \approx 0.46999\). مع مطابقة مقياس بلانك \(\LUV = 3\pi\Mpl\)، يكون التصحيح الإشعاعي لكتلة هيغز \(\delta m_H \approx 1.57\times 10^{18}\) GeV، مما يتطلب إلغاء مضبوطاً بدقة مع الكتلة العارية.
		
		كبديل، اقترحنا السيناريو B، حيث تكون لاغرانجية العارية على مقياس بلانك ثابتة مقياسياً وتتولد جميع الكتل إشعاعياً. عندئذٍ يطلب الاتساق الذاتي أن يكون المقياس الأساسي للتنسيق \(\LUV \approx 8.96\) TeV – قيمة يمكن الوصول إليها في LHC. يقدم هذا السيناريو تنبؤات ملموسة قابلة للتزييف لعمليات الاستطارة عالية الطاقة.
		
		لا يمكن اتخاذ الخيار بين السيناريو A (غير مرضي جمالياً لكنه متسق مع أصل مقياس بلانك لـ YUCT) والسيناريو B (تنبؤي وقابل للاختبار، لكنه يتطلب إعادة تفسير سلم كتلة YUCT) على أسس نظرية وحدها. يجب أن يقرره التجربة. نحث تعاونيات LHC على فحص بيانات دريل-يان والنفاثات المزدوجة عالية الكتلة بحثاً عن أدلة على كبح عامل الشكل الذي تنبؤ به المعادلة (\ref{eq:F}) مع \(\LUV \sim 9\) TeV. النتيجة الصفرية تستبعد السيناريو B، وبينما يبقى السيناريو A ممكناً منطقياً، فإنه يضعف قضية YUCT كحل لمشكلة التسلسل الهرمي.
		
		\section*{شكر وتقدير}
		يشكر المؤلف المشاركين في ندوة YUCT على المناقشات النقدية.
		
		\section*{توفر البيانات}
		جميع الحسابات مكتفية ذاتياً؛ كود بايثون المستخدم للتقييم العددي متاح في مستودع YPSDC: \url{https://github.com/Alexey-Yakushev-YUCT/YPSDC}.
		
		\begin{thebibliography}{99}
			\bibitem{AppendixL} A.\,V.\,Yakushev, “الملحق L: القياس الكسيري لخطأ التنسيق,” في \emph{YUCT}, Zenodo (2026).
			\bibitem{AppendixV} A.\,V.\,Yakushev, “الملحق V: الزمن كإنتروبيا عمليات التنسيق,” في \emph{YUCT}, Zenodo (2026).
			\bibitem{AppendixG} A.\,V.\,Yakushev, “الملحق G: الجاذبية كشد هندسي لشبكة التنسيق وعلم الكون الدوري في YUCT,” في \emph{YUCT}, Zenodo (2026).
			\bibitem{AppendixM} A.\,V.\,Yakushev, “الملحق M: علم كون YUCT – التضخم كتحول طوري تنسيقي,” في \emph{YUCT}, Zenodo (2026).
			\bibitem{AppendixLambda} A.\,V.\,Yakushev, “الملحق \(\Lambda\): حل مشكلتي التسلسل الهرمي والثابت الكوني في YUCT,” في \emph{YUCT}, Zenodo (2026).
			\bibitem{AppendixJ} A.\,V.\,Yakushev, “الملحق J: اشتقاق كامل لبارامترات النكهة في النموذج المعياري من الإزاحات الطوبولوجية لـ YUCT,” في \emph{YUCT}, Zenodo (2026).
			\bibitem{AppendixFerra} A.\,V.\,Yakushev, “الملحق Ferra1.2: ثوابت التنسيق العامة للأطوار المنتظمة وغير المنتظمة,” في \emph{YUCT}, Zenodo (2026).
		\end{thebibliography}
		
	\end{otherlanguage}
	
\end{document}